\selectlanguage{ngerman}%
\frenchspacing%
\begin{ledgroupsized}[r]{114mm}%
\footnotesize%
\pstart%
\noindent%
\textbf{Überlieferung:}%
\pend%
\end{ledgroupsized}%
\begin{ledgroupsized}[r]{114mm}%
\footnotesize%
\pstart%
\parindent -6mm%
\makebox[6mm][l]{\textit{E}}%
Druck nach unbekannter Vorlage:
\textsc{Gerland} 1906, S.\;9\,f.
(Nr.\;6 mit der Überschrift \textit{Kapillarität}).%
\cite{00197}
Die verschollene handschriftliche Vorlage wies laut \textit{E} (S.\;9) das Format eines kleinen Blättchens auf.
Der oft fehlerhafte Text von \textit{E} (Gerlands Transkription) wird hier wiedergegeben und gegebenenfalls berichtigt oder kommentiert;
die Rechtschreibung wird an die Editionsregeln angepasst.
\newline%
Cc~2 Nr.\;0
\pend%
\end{ledgroupsized}%
\vspace{5mm}%
\begin{ledgroup}%
\footnotesize%
\pstart%
\noindent%
\textbf{Datierungsgründe:}
\pagelabel{56216_dat}%
Die vorliegende Aufzeichnung N.\;\ref{56216}, deren handschriftliche Vorlage nach heutigem Wissensstand verschollen ist, ist undatiert.
Zu Beginn des in \textit{E} posthum überlieferten Textes bezieht sich Leibniz aber eindeutig auf die 1674 in Paris veröffentlichte Abhandlung \textit{De l'origine des fontaines}\cite{02534} von Pierre Perrault\indexN{\textls{Perrault}, Pierre 1611\textendash1680}.
Die Genauigkeit der Bezugnahme zeigt, dass Leibniz das Buch wohl aus erster Hand gekannt haben muss.
Dies lässt wiederum vermuten, dass er es noch vor seiner Abreise aus Paris\indexO{Paris} Anfang Oktober 1676 gelesen hatte.
Laut abgedrucktem königlichem Privileg wurde die Abhandlung erst am 18.\;September 1674 vom Drucker ausgeliefert (vgl. \textsc{P.\;Perrault}‚ \textit{De l'origine des fontaines}, Paris 1674, S.\,e\,ii\,v\textsuperscript{o}).\cite{02534}
Folglich kann Leibniz sie wohl frühestens im Oktober 1674 gelesen haben.
Dass seine Lektüre von Perraults Abhandlung erst nach der Abreise aus Paris\indexO{Paris} erfolgte, ist zwar möglich, aber weniger wahrschenlich.
Als Datierung der vorliegenden Aufzeichnung N.\;\ref{56216} wird demgemäß die Zeitspanne vom Oktober 1674 bis zum September 1676 (unter Vorbehalt) angenommen.
\pend%
\end{ledgroup}%
\selectlanguage{latin}%
\frenchspacing%
\vspace{8mm}%
\normalsize%
\pstart%
\noindent%
\bracks{S.\,9}
Es ist gewiß,
daß das waßer%
\indexS{Waßer}
über seine Wage%
\indexS{Waage}
in die höhe steiget,
sowohl in engen röhren,%
\indexS{Rohre}
als in sand %
\indexS{Sand}
schwamm%
\indexS{Schwamm}
und dergleichen%
\bracks{,}
an die 18 Zoll%
\indexS{Zoll}%
\bracks{:}
videatur
\edtext{%
Experimentum Maignani%
\indexN{\textls{Maignan} (Magnanus), Emmanuel 1633\textendash1676}%
\indexS{experimentum Maignani}
relatum et excultum a Dn
\edtext{%
\bracks{Perr}alto%
\indexN{\textls{Rohault}, Jacques 1618\textendash1672}%
\indexN{\textls{Perrault}, Pierre 1611\textendash1680}%
}{%
\lemma{Rohalto}%
\Bfootnote{%
\textit{E\;ändert Hrsg.}%
}}
libro
\edtext{%
\bracks{\textit{de}}%
}{%
\lemma{die}%
\Bfootnote{%
\textit{E\;ändert Hrsg.}%
}}
\textit{origine fontium},%
}{%
\lemma{Experimentum \bracks{...} \textit{fontium}}%
\Cfootnote{%
Siehe
\textsc{P.\;Perrault},
\textit{De l'origine des fontaines},
partie II
(Paris 1674, S.\;154--160).%
\cite{02534}
Berichtet über das Experiment hatte ursprünglich
\textsc{E.\;Maignan},
\textit{Cursus philosophicus} %
(Toulouse 1652; 2.\;Ausgabe, Lyon 1673).%
\cite{02535}
Das Referat hierüber in
\textsc{K.\;Schott},\indexN{\textls{Schott} (Schottus), Kaspar 1608\textendash1666}
\textit{Anatomia physico-hydrostatica fontium ac fluminum},
lib.\;IV, cap.\;14, §\;4
(Würzburg 1663, S.\;254\,f.)%
\cite{02535}
könnte aber Perraults\indexN{\textls{Perrault}, Pierre 1611\textendash1680} unmittelbare Quelle gewesen sein;
vgl. %
\textit{De l'origine des fontaines},
partie I
(S.\;125\,f.).%
\cite{02534}%
}}
sed ut effluat seu redundet ultra tubulos,%
\indexS{tubulus angustus}
in quibus ascendit,
effici non potest.
Ego vero pro rectum habeo
non id
\bracks{S.\,10}
fieri
\edtext{%
ex propria vi aequilibrii,%
\indexS{vis aequilibrii}%
\indexS{aequilibrium}
ut concipere videtur Robervallus,%
\indexN{\textls{Roberval} (Robervallius), Gilles Personne de 1602\textendash1675}
quod aqua%
\indexS{aqua aequiponderans}
in tubulis angustis%
\indexS{tubulus angustus}
non aequiponderet,%
}{%
\lemma{ex propria vi \bracks{...} aequiponderet}%
\Cfootnote{%
Nicht ermittelt.%
}}
sed
\edtext{%
ex vi externa%
\indexS{vis externa}
circiter,
ut explicat Rohaltus.%
\indexN{\textls{Rohault}, Jacques 1618\textendash1672}%
}{%
\lemma{ex vi externa \bracks{...} Rohaltus}%
\Cfootnote{%
Siehe
\textsc{J.\;Rohault},
\textit{Traité de physique},
partie~I, cap.\;XXII, §\;80--89
(Paris 1671, S.\;190--195).%
\cite{00087}%
}}
Quoniam igitur adest illa vis,%
\indexS{vis externa}
videndum est,
an non revera eius ope effici possit motus perennis,%
\indexS{motus perennis mechanicus}
\ledpb%
qui tamen revera non foret mechanicus%
\indexS{motus perennis mechanicus}%
\bracks{,}
non magis %
quam
\edtext{%
ille %
quem
\edtext{%
Leisnerius%
\indexN{\textls{Lersner}, Johann Maximilian, gest. 1702}%
}{%
\lemma{Leisnerius}%
\Cfootnote{%
Möglicherweise Johann Maximilian Lersner (gest.\;1702), der in den Jahren 1671--1672 mit Leibniz korrespondierte.
Die Form \textit{Leisnerius} könnte von einem Transkriptionsfehler herrühren.%
}}
ope Magnetis%
\indexS{motus perennis ope magnetis}%
\indexS{magnes}
procurare voluerat.%
}{%
\lemma{ille \bracks{...} voluerat}%
\Cfootnote{%
Nicht ermittelt.%
}}
Adhibeantur tenuissima folia%
\indexS{folium talci tenue}
ut talci%
\indexS{talcum}%
\bracks{,}
et eorum motu,
dum rursus a se dimoventur,
\edtext{%
recidat%
\bracks{,}%
}{%
\lemma{recidat}%
\Cfootnote{%
Die Pluralform \textit{recidant} wäre angemessener;
die Bedeutung des Satzes bleibt jedoch dunkel.
Hier ist der in \textit{E} überlieferte Text wohl verderbt.%
}}
vel si circuletur aliquid
perpendiculariter erectum
circa subtilissimos axes,%
\indexS{axis}
ut eius motu modo fiant,
modo
\edtext{%
destin\bracks{e}ntur%
}{%
\lemma{destinantur}%
\Bfootnote{%
\textit{L\;ändert Hrsg.}%
}%
\lemma{destin\bracks{e}ntur}%
\Cfootnote{%
Die richtige Lesart könnte \textit{desinentur} sein.%
}%
}
subtiles crenae.%
\indexS{crena}
Experiamur etiam,
quod fiat,
si
\edtext{%
quam continua pagina venae~(?)%
}{%
\lemma{quam \bracks{...} venae~(?)}%
\Cfootnote{%
Von \textit{E} selbst als unsicher ausgewiesene Lesung.%
}}
alicubi angustior alicubi laxior sit%
\bracks{;}
erunt enim et liquores%
\indexS{liquores altitudinis inaequalis}
\edtext{%
inaequal\bracks{i}s%
}{%
\lemma{inaequales}%
\Bfootnote{%
\textit{E\;ändert Hrsg.}%
}}
altitudinis.
Quid si foliolis%
\indexS{foliolum}
ab intus accedat exigua pressio,%
\indexS{pressio exigua}
ut liquorem exprimat%
\indexS{liquor expressus}%
\bracks{:}
an sponte alius succedet,
fluentque etiam recessante pressione?
Non puto,
alioqui res foret egregia:%
\indexS{res egregia}
puto igitur potius
subinde repetita pressione%
\indexS{pressio repetita}
opus fore.
Videndum autem,
an vis exprimens%
\indexS{vis exprimens}
minor sufficiat,%
\indexS{vis sufficiens}
quam est ipsum aquae elevatae pondus.%
\indexS{pondus aquae elevatae}%
\indexS{aqua elevata}
Certe non solum pondus aquae superandum est,%
\indexS{pondus aquae superandum}
sed et adhaesio ad superficiem%
\indexS{adhaesio ad superficiem}
foliolorum.%
\indexS{foliolum}
\pend%
\selectlanguage{latin}%
\frenchspacing%
\normalsize%
